\subsection{李群芽的李代数}

本号中，\((G, e, \theta, m)\) 表示一个李群芽。本节的大部分结果仍以同样的证明成立。我们复述其中将要用到的结果。

18.1. 设 \(\Omega\) 是 \(m\) 的定义域。设 \((g, g') \in \Omega\)、\(t \in \mathrm{T}_g^{(\alpha)}(\mathrm{G})\)、\(t' \in \mathrm{T}_{g'}^{(\alpha)}(\mathrm{G})\)。如第1号所述，\(t\) 与 \(t'\) 的卷积乘积记为 \(t * t'\)，它是 \(t \otimes t'\) 在 \(m\) 下的像。记 \(\mathrm{U}(\mathrm{G}) = \mathrm{T}_e^{(\alpha)}(\mathrm{G})\)、\(\mathrm{U}_s(\mathrm{G}) = \mathrm{T}_e^{(s)}(\mathrm{G})\)、\(\mathrm{U}^+(\mathrm{G}) = \mathrm{T}_e^{(\alpha)+}(\mathrm{G})\)、\(\mathrm{U}_s^+(\mathrm{G}) = \mathrm{T}_e^{(s)+}(\mathrm{G})\)。对于 \(t\) 和 \(t'\)（它们属于 \(\mathrm{U}(\mathrm{G})\)），\(t * t'\) 有定义且属于 \(\mathrm{U}(\mathrm{G})\)。在卷积乘积下，\(\mathrm{U}(\mathrm{G})\) 是单位元为 \(\varepsilon_e\) 的结合代数，并由 \(\mathrm{U}_s(\mathrm{G})\) 滤化。典范同构 \(i_{\mathrm{G}, e}\) 将 gr \(\mathrm{U}(\mathrm{G})\) 映到 \(\mathrm{TS}(\mathrm{T}_e(\mathrm{G}))\) 上，并且是代数同构。

18.2. 设 G、H 是李群芽，\(\phi: G \to H\) 是一个态射。若 \(t \in \mathrm{U}(\mathrm{G})\)，则 \(\mathrm{U}(\phi)(t)\) 是 t 在 \(\phi_{*}\) 下的像，是 U(H) 的元素，且 \(\mathrm{U}(\phi)\) 是从代数 U(G) 到代数 U(H) 的态射。G 到 G 的映射 \(\theta: x \mapsto x^{-1}\) 定义了 U(G) 到 U(G) 的映射 \(t \mapsto t^{\vee}\)。对于 U(G) 中的 \(t, t'\)，\(t * t'\) 相对于 \(G^{\vee}\) 计算时，等于 \(t' * t\) 相对于 G 计算时的结果，且 \((t * t')^{\vee} = t^{\prime\vee} * t^\vee\)。于是 \(\mathrm{U}(\phi)(t^{\vee}) = (\mathrm{U}(\phi)t)^{\vee}\)。若 \(G_{1}, \ldots, G_{n}\) 是李群芽且 \(G = G_{1} \times \cdots \times G_{n}\)，则从 \(\mathrm{U}(G_{1}) \otimes \cdots \otimes \mathrm{U}(G_{n})\) 到 U(G) 的典范同构是代数同构；

对 U(G) 中的 \(t_{1},\ldots,t_{n}\)，\((t_{1}\otimes\cdots\otimes t_{n})^{\vee}=t_{1}^{\vee}\otimes\cdots\otimes t_{n}^{\vee}\)。设 H 是 G 的李子群芽，i: H→G 是典范嵌入。则 U(i) 是从代数 U(H) 到 U(G) 的单射同态，并且

\[
\mathrm{U} (i) \left(t ^ {\vee}\right) = \left(\mathrm{U} (i) (t)\right) ^ {\vee}
\]

对所有 \(t \in \mathrm{U}(\mathrm{H})\) 均成立。在由 G 的流形结构定义的卷积乘积和余乘法下，\(\mathrm{U}(\mathrm{G})\) 是双代数，\(\mathrm{U}(\phi)\) 是双代数态射。

18.3. 设 G 是李群芽，X 是 \(C^{r}\) 类流形，\(\psi\) 是 G 在 X 上的 \(C^{r}\) 类左作用律片段。设 \(\Omega\) 是 \(\psi\) 的定义域。若 \(t \in T_{g}^{(s)}(G)\)、\(u \in T_{x}^{(s^{\prime\prime})}(X)\)、\((g, x) \in \Omega\) 且 \(s + s^{\prime} \leqslant r\)，则令 \(t * u\) 表示 \(t \otimes u\) 在 \(\psi_{*}\) 下的像。设 \(t \in T_{g}^{(s)}(G)\)、\(t^{\prime} \in T_{g'}^{(s^{\prime})}(G)\)、\(u \in T_{x}^{(s^{\prime})}(X)\)；若 \(s + s^{\prime} + s^{\prime\prime} \leqslant r\) 且 \(gg^{\prime}\)、\((gg^{\prime})x\)、\(g^{\prime}x\)、\(g(g^{\prime}x)\) 有定义，则

\[
(t * t ^ {\prime}) * u = t * (t ^ {\prime} * u).
\]

设 \(x_0 \in \mathbf{X}\)，\(\rho(x_0)\) 是映射 \(g \mapsto gx_0\)，它在 \(e\) 的某个开邻域上有定义。若 \(t \in \mathrm{U}_r(\mathrm{G})\)，则 \(\rho(x_0)_*t = t * \varepsilon_{x_0}\)。这里以及本号其余部分，右作用律片段的相应结果留给读者翻译。

18.4. 保持 18.3 的记号，设 \(t \in \mathrm{U}_s(\mathbf{G})\) 且 \(s \leqslant r\)。设 \(f\) 是 \(\mathbf{C}^r\) 类函数，定义在 \(\mathbf{X}\) 上且取值于 Hausdorff 多赋范空间。令 \(t * f\) 表示由下式定义的 \(\mathbf{X}\) 上的函数：

\[
\begin{array}{r l} (t * f) (x) & = \langle t, g \mapsto f (\psi (\theta (g), x)) \rangle \\ & = \langle t ^ {\vee}, f \circ \rho (x) \rangle = \langle \rho (x) _ {*} (t ^ {\vee}), f \rangle = \langle t ^ {\vee} * \varepsilon_ {x}, f \rangle . \end{array}
\]

若 \(t \in \mathrm{U}_s(\mathrm{G})\)、\(t' \in \mathrm{U}_{s'}(\mathrm{G})\) 且 \(s + s' \leqslant r\)，则 \(\langle t', t * f \rangle = \langle t^\vee * t', f \rangle\) 且 \((t * t') * f = t * (t' * f)\)。设 \(t \in \mathrm{U}_s(\mathrm{G})\)，\(f\) 和 \(f'\) 是 \(C^r\) 类函数，定义在 X 上且分别取值于 Hausdorff 多赋范空间 \(F, F'\)，\((u, u') \mapsto u. u'\) 是从 \(F \times F'\) 到 Hausdorff 多赋范空间的连续双线性映射；令 \(\underline{n}\)。

\(\sum_{i=1}^{n} t_i \otimes t_i'\) 是 \(t\) 在余乘法下的像；若 \(s \leqslant r\)，则

\[
t * (f f ^ {\prime}) = \sum_ {i = 1} ^ {n} (t _ {i} * f) (t _ {i} ^ {\prime} * f ^ {\prime}).
\]

18.5. 保持 18.3 的记号，设 \(t \in \mathrm{U}_s(\mathrm{G})\) 且 \(s \leqslant r\)。映射 \(x \to t * \varepsilon_x\) 称为由 \(t\) 和作用律片段定义的点分布场，有时记为 \(\mathrm{D}_t^{\psi}\) 或 \(\mathrm{D}_t\)。若 \(f: \mathrm{X} \to \mathrm{F}\) 是 \(C^r\) 类函数，则函数 \(t^{\vee} * f\) 在 \(\mathrm{X}\) 上也记为 \(\mathrm{D}_t f\)；它属于 \(C^{r - s}\) 类（当 \(s < \infty\) 时）。若 \(t \in \mathrm{U}_s(\mathrm{G})\)、\(t' \in \mathrm{U}_{s'}(\mathrm{G})\) 且 \(s + s' \leqslant r\)，则 \(\mathrm{D}_{t,t'} f = \mathrm{D}_{t'}(\mathrm{D}_t f)\)。若 \(G\) 和 \(X\) 是有限维的，则 \(D_t\) 是 \(X\) 上阶数 \(\leqslant s\) 且属于 \(C^{r - s}\) 类的微分算子（若 \(s < \infty\)）。函数 \(D_t f\) 就是 \(f\) 在这个微分算子下的变换。

18.6. 设 G 是李群芽且 \(t \in \mathrm{U}(\mathrm{G})\)。\(L_{t}\) 表示 G 上的点分布场 \(g \mapsto \varepsilon_{g} * t\)，而 \(R_{t}\) 表示点分布场 \(g \mapsto t * \varepsilon_{g}\)。若 \(f \in \mathcal{C}^{\omega}(\mathrm{G}, \mathrm{F})\)，则 \(\mathrm{L}_{t} f \in \mathcal{C}^{\omega}(\mathrm{G}, \mathrm{F})\) 且 \(\mathrm{R}_{t} f \in \mathcal{C}^{\omega}(\mathrm{G}, \mathrm{F})\)。对于 \(t, t'\)，它们属于 \(\mathrm{U}(\mathrm{G})\)，\(L_{t* t'} = L_t \circ L_{t'}, R_{t* t'} = R_{t'} \circ R_t, L_t \circ R_{t'} = R_{t'} \circ L_t, \theta(L_t) = R_t^\vee\)。

18.7. 由于 \(T_{e}(G)\) 是 U(G) 的本原元素集合，

\[
[ \mathrm{T} _ {e} (\mathrm{G}), \mathrm{T} _ {e} (\mathrm{G}) ] \subset \mathrm{T} _ {e} (\mathrm{G}).
\]

可赋范空间 \(\mathrm{T}_{e}(\mathrm{G})\) 与括号一起构成可赋范李代数，称为 G 的可赋范李代数（或 G 的李代数），记为 \(\mathrm{L}(\mathrm{G})\)。设 \(\mathrm{E}(\mathrm{G})\) 是 \(\mathrm{L}(\mathrm{G})\) 的包络代数。\(\mathrm{L}(\mathrm{G})\) 到 \(\mathrm{U}(\mathrm{G})\) 的典范嵌入定义了同态 \(\eta\)，它从代数 \(\mathrm{E}(\mathrm{G})\) 到代数 \(\mathrm{U}(\mathrm{G})\)；若 K 的特征为 0，则 \(\eta\) 是双代数同构，借此将 \(\mathrm{U}(\mathrm{G})\) 与 \(\mathrm{E}(\mathrm{G})\) 识别。采用 18.3 的记号，对所有 \(a \in \mathrm{L}(\mathrm{G})\)，设 \(D_{a}\) 是由 a 在 X 上定义的点分布场。映射 \((a, x) \mapsto D_{a}(x)\) 是 \(C^{r-1}\) 类态射，从平凡向量丛 \(\mathrm{L}(\mathrm{G}) \times \mathrm{X}\) 到向量丛 \(\mathrm{T}(\mathrm{X})\)。设 I 是包含 0 的 K 的一个开子集，\(\gamma: I \to G\) 是 \(C^{r}\) 类映射，满足 \(\gamma(0) = e\)。令 \(a = T_{0}(\gamma) 1 \in L(\mathrm{G})\)。若 f: X → F 是 \(C^{r}\) 类函数，则

\[
\left(\mathrm{D} _ {a} f\right) (x) = \lim _ {k \in \mathrm{K} ^ {*}, k \rightarrow 0} k ^ {- 1} \left(f (\gamma (k) x) - f (x)\right).
\]

若 \(r \geqslant 2\)，则映射 \(a \mapsto \mathrm{D}_a\) 是 \(\mathrm{C}^{r-1}\) 类的 \(\mathrm{L}(\mathrm{G})\) 在 \(\mathrm{X}\) 上的左无穷小作用律。

18.8. 设 G 和 H 是李群芽，\(\phi\) 是从 G 到 H 的态射。\(\mathrm{U}(\phi)\) 在 \(\mathrm{L}(\mathrm{G})\) 上的限制，也就是 \(\mathrm{T}_{e}(\phi)\)，是从 \(\mathrm{L}(\mathrm{G})\) 到 \(\mathrm{L}(\mathrm{H})\) 的连续态射，记为 \(\mathrm{L}(\phi)\)。若 \(\psi\) 是从 H 到某个李群芽的态射，则 \(\mathrm{L}(\psi \circ \phi) = \mathrm{L}(\psi) \circ \mathrm{L}(\phi)\)。\(\phi\) 是浸入的充要条件是，\(\mathrm{L}(\phi)\) 是从 \(\mathrm{L}(\mathrm{G})\) 到 \(\mathrm{L}(\mathrm{H})\) 中某个具有拓扑补空间的李子代数的同构。特别地，若 G 是 H 的李子群芽且 \(\phi\) 是典范嵌入，则将 \(\mathrm{L}(\mathrm{G})\) 与 \(\mathrm{L}(\mathrm{H})\) 的一个李子代数借助 \(\mathrm{L}(\phi)\) 识别。若 \((\mathrm{G}_{i})_{i\in I}\) 是李群芽的有限族且 G 是它们的乘积，则 \(\mathrm{L}(\mathrm{G})\) 与 \(\prod_{i\in I}\mathrm{L}(\mathrm{G}_{i})\) 典范地识别。

18.9. 设 G 是李群芽（当 K 的特征 \textgreater0 时为有限维）。设 F 是完备可赋范空间。设 \(\alpha\) 是 G 上取值于 F 的 k 阶微分形式。称 \(\alpha\) 在 G 上左不变，当 \(\alpha_{g}\) 由 \(\alpha_{e}\) 经映射 \(h \mapsto gh\) 从 e 的某个邻域导出到 g 的某个邻域时。若 \(\alpha\) 左不变，则 \(\alpha\) 解析。映射 \(\alpha \mapsto \alpha_{e}\) 是从 G 上取值于 F 的 k 阶左不变微分形式集合到从 \(T_{e}(G)\) 到 F 的连续交替 k-线性映射集合的双射。若 \(\alpha_{e} = \operatorname{Id}_{\mathrm{T}_e(\mathrm{G})}\)，则称 \(\alpha\) 为 G 的左典范微分形式。右不变微分形式及右典范微分形式的定义类似。若 \(\omega\) 是 G 的左典范微分形式，则 \(d\omega + [\omega]^{2} = 0\)。设 M 是 \(C^{r}\) 类流形，f 是从 M 到 G 的 \(C^{r}\) 类映射。f 的左微分记为 \(f^{-1}.df\)，它是 M 上取值于 L(G) 的 1 阶微分形式，把每个向量 \(u \in T_{m}(M)\) 对应到元素 \(f(m)^{-1}\) . \((\mathrm{T}_{m}f)(u)\)。则 \(f^{-1}.df = f^{*}(\omega)\) 且 \(d\alpha + [\alpha]^{2} = 0\)。若从 M 到 G 的两个映射 f 和 g 具有相同的左微分且 K 的特征为 0，则 \(fg^{-1}\) 局部为常值。

